\section{Original analytic prerequisites and submanifold charts}

\begin{lemma}[A closed complex submanifold is a complex manifold, and its inclusion an immersion]
\label{ana:submanifold-complex-structure}
\uses{ana:geometric-setting,ana:blowups}
\lean{Hironaka.SubmanifoldChartedSpace.chartedSpace,Hironaka.SubmanifoldChartedSpace.isComplexManifold_chartedSpace,Hironaka.SubmanifoldChartedSpace.contMDiff_val}\upstreamlean{TranscendentalNakamaye.Pending.StepFourOnStrictTransform.fderiv_coordinateMap_val_injective}
\mixedleanok
Let \(W\) be a subset of a complex manifold \(X\) of dimension \(n\) which is regular of
dimension \(k\le n\) at each of its points. With the subspace topology and the atlas of
flattening charts (at each point, a holomorphic chart of \(X\) followed by a biholomorphism
straightening \(W\) to \(\mathbf C^k\times0\), restricted to \(W\)), \(W\) is a complex manifold
of dimension \(k\), and the inclusion \(W\hookrightarrow X\) is holomorphic, with injective
differential in every pair of charts.
\end{lemma}
\begin{proof}
\mixedleanok
The transition between two flattening charts is the composite of holomorphic maps restricted
to \(\mathbf C^k\times0\), hence holomorphic; Hausdorffness and second countability pass to
the subspace. In a flattening chart \(F\) and the ambient chart it is built from, the inclusion
reads \(\zeta\mapsto\Phi_F^{-1}(\zeta,0)\), which is holomorphic. In a flattening chart \(F\) and an
ambient chart \(d\), the inclusion read in coordinates has the holomorphic left inverse \(w\mapsto\operatorname{pr}_k(F(d^{-1}(w)))\)
near the image point, so its differential is injective.
\end{proof}

\begin{lemma}[Holomorphic maps into a submanifold]
\label{ana:submanifold-corestriction}
\uses{ana:submanifold-complex-structure}
\source{transcendental-nakamaye-copy4-20260911:prop:principalization}
\lean{Hironaka.SubmanifoldCodRestrict.contMDiff_codRestrict,Hironaka.SubmanifoldCodRestrict.contMDiff_restrictPreimage}
\leanok
Let \(W\subseteq X\) be as in Lemma~\ref{ana:submanifold-complex-structure}, with its complex
structure. If \(f:M\to X\) is a holomorphic map of complex manifolds with \(f(M)\subseteq W\), then
\(f\), regarded as a map \(M\to W\), is holomorphic. In particular, if \(\pi:Y\to X\) is holomorphic
and \(W'\subseteq Y\) is regular of some dimension \(k'\) at each of its points, with
\(\pi(W')\subseteq W\), then the restriction \(\pi|_{W'}:W'\to W\) is holomorphic.
\end{lemma}
\begin{proof}
\leanok
Let \(y\in M\) and let \(F\) be the flattening chart of \(W\) at \(f(y)\): a chart \(e_F\) of \(X\)
followed by a biholomorphism \(\Phi_F\) straightening \(W\) to \(\mathbf C^k\times0\). The map
\(f\) is continuous into \(W\) for the subspace topology, and in the chart \(F\) it reads
\(z\mapsto\operatorname{pr}_k(\Phi_F(e_F(f(z))))\) near \(y\), a composite of holomorphic maps. For the
second assertion apply the first to \(\pi\circ\iota_{W'}\), where the inclusion
\(\iota_{W'}:W'\to Y\) is holomorphic by Lemma~\ref{ana:submanifold-complex-structure}.
\end{proof}

\begin{lemma}[Blowup sequences are surjective]
\label{ana:blowup-surjective}
\uses{ana:blowups,ana:strict-transform-composition}
\lean{Hironaka.StrictTransformComposition.surjective_of_isBlowupAlong,Hironaka.StrictTransformComposition.IsBlowupSequence.compactSpace_surjective,Hironaka.StrictTransformComposition.image_exceptionalLocus_subset_comp,Hironaka.StrictTransformComposition.not_strictTransform_subset_image_exceptionalLocus}
\leanok
Let \(X\) be a compact complex manifold and \(\mu:\widetilde X\to X\) a finite
composition of blowups with smooth centres. Then \(\widetilde X\) is compact and
\(\mu\) is surjective. If \(f:Y\to X\) and \(g:Z\to Y\) are surjective
holomorphic maps between manifolds of the same dimension and
\(V\not\subset(f\circ g)(\operatorname{Exc}(f\circ g))\), then
\(V\not\subset f(\operatorname{Exc}(f))\) and
\(\widetilde V^{f}\not\subset g(\operatorname{Exc}(g))\).
\end{lemma}
\begin{proof}
\leanok
A blowup \(g:Z\to Y\) along a nowhere dense centre \(C\) of a compact manifold is
an isomorphism over \(Y\setminus C\), and \(Z\) is compact. Its image is compact,
hence closed, and contains the dense set \(Y\setminus C\), so \(g\) is onto. By
induction along the sequence every intermediate manifold is compact and the
composition is onto.

If \(y=f(x)\) with \(x\in\operatorname{Exc}(f)\), write \(x=g(w)\); then
\(w\in g^{-1}\operatorname{Exc}(f)\subset\operatorname{Exc}(f\circ g)\), so
\(f(\operatorname{Exc} f)\subset(f\circ g)(\operatorname{Exc}(f\circ g))\), which
gives the first claim. For the second, suppose
\(\widetilde V^{f}\subset g(\operatorname{Exc} g)\) and let \(v\in V\), \(v=f(y)\).
If \(y\in\operatorname{Exc}(f)\), then \(v\in(f\circ g)(\operatorname{Exc}(f\circ g))\)
as above. Otherwise \(y\in f^{-1}(V)\setminus\operatorname{Exc}(f)\subset\widetilde V^{f}\),
so \(y=g(w)\) with \(w\in\operatorname{Exc}(g)\subset\operatorname{Exc}(f\circ g)\)
and \(v=(f\circ g)(w)\). Thus \(V\subset(f\circ g)(\operatorname{Exc}(f\circ g))\), a
contradiction.
\end{proof}

\begin{proposition}[A connected manifold is irreducible]
\label{ana:manifold-irreducible}
\uses{fpgeom:complex-analytic-subvariety,ana:nonkahler-locus,ana:blowups}
\source{demailly-cadg:II.4.2}
\lean{Hironaka.AnalyticIrreducibleComponents.isComplexAnalyticallyIrreducible_univ,Hironaka.AnalyticIrreducibleComponents.regularAtDim_univ}\upstreamlean{TranscendentalNakamaye.Pending.AnalyticIrreducibleComponents.exists_subvariety_univ,TranscendentalNakamaye.Pending.AnalyticIrreducibleComponents.eq_of_regularAtDim_univ,TranscendentalNakamaye.Pending.AnalyticIrreducibleComponents.dim_eq_of_carrier_eq_univ}\lean{Hironaka.AnalyticIrreducibleComponents.isOpen_of_forall_regularAtDim_self}\upstreamlean{TranscendentalNakamaye.Pending.AnalyticIrreducibleComponents.eq_univ_of_isClosedConnectedSubmanifold_self}
\mixedleanok
Let \(X\) be a connected complex manifold of dimension \(n\).
\begin{enumerate}
\item \(X\) is an irreducible analytic subset of itself, of dimension \(n\), and
      \(n\) is the only dimension of a regular point of \(X\).
\item A closed connected complex submanifold of \(X\) of dimension \(n\) is \(X\).
\end{enumerate}
\end{proposition}
\begin{proof}
\mixedleanok
(1) Every point of \(X\) is a regular point of dimension \(n\). If
\(X=B\cup C\) with \(B,C\) analytic and \(B\ne X\), then \(B\) has empty
interior: an analytic subset containing a nonempty open set contains, by the
identity theorem applied to its local equations along chains of coordinate balls,
the whole connected manifold. Hence the closed set \(C\) contains the dense set
\(X\setminus B\), so \(C=X\). A flattening of \(X\) of dimension \(k\) at a point
maps an open neighbourhood of that point onto an open set contained in
\(\mathbb C^k\times\{0\}\); a nonempty open subset of \(\mathbb C^n\) lies in that
plane only if \(k=n\).

(2) At each of its points a submanifold of dimension \(n\) contains a
neighbourhood of that point, so it is open. It is also closed and nonempty, hence
equal to the connected manifold \(X\).
\end{proof}
